\lecture{5}{Mon 12 Sep 2022 10:30}{}
\begin{nonexample}
 	Compare $\left( \mathbb{Z}_6, + \right) $ and $D_6 = \{\text{id}, \sigma, \sigma^2, \tau,\tau\sigma, \tau\sigma^2\} $. They both have 6 elements, hence there exists a bijection from $\mathbb{Z}_6$ onto $D_6$. However there is no isomorphism.
\end{nonexample}

\begin{example}
	\begin{enumerate}
		\item (Trivial homomorphism) \begin{align*}
			\varphi: G &\longrightarrow G' \\
			g &\longmapsto \varphi(g) = e'
		.\end{align*}
	\item Consider \begin{align*}
		\varphi: \mathbb{Z} &\longrightarrow \mathbb{Z}_n \\
		a &\longmapsto \varphi(a) = [a] = a \mod n
	.\end{align*}
	Then $\varphi(a+b) = [a+b]=[a]+[b] = \varphi(a)+\varphi(b)$.
	\item (Identity map is a trivial homomorphism) 
		\begin{align*}
			\varphi: G &\longrightarrow G \\
			g &\longmapsto \varphi(g) = g
		.\end{align*}
	\item \begin{align*}
		\varphi: \left( \mathbb{Z}, + \right)  &\longrightarrow \left( \mathbb{Q}^\times , \times  \right)  \\
		n &\longmapsto \varphi(n) = 2^n
	.\end{align*}
	Now $\varphi(a+b) = 2^{a+b} = 2^a 2^b = \varphi(a)\varphi(b)$.

	\item \begin{align*}
		\varphi: \left( \mathbb{Q}^\times , \times  \right)  &\longrightarrow \left( \mathbb{Q}^\times , \times  \right)  \\
		x &\longmapsto \varphi(x) = x^2
	.\end{align*}
	Now $\varphi(a+b) = (ab)^2 = a^2b^2 = \varphi(a)\varphi(b)$.
	\end{enumerate}
\end{example}

\begin{theorem}[Cayley's Theorem]
	Every group $G$ is isomorphic to a subgroup of $A(S)$ for some set $S$.
\end{theorem}
\begin{corollary}
	When $G$ is a finite group, then $G$ is isomorphic to a subgroup of $S_n$ for some $n$.
\end{corollary}

\begin{proof}
	Let $G$ be a group, and put $S = G$. We define the set $A(G)$ to be the set of bijective self-mapping of $G$. Observe that whenever $a \in G$, the map \begin{align*}
		T_a: G &\longrightarrow G \\
		g &\longmapsto T_a(g) = ag
	.\end{align*}
	belongs to $A(G)$. To check bijectivity, observe that $T_a^{-1} := T_{a^{-1}}$ (for all $a$, $(T_a^{-1} \circ T_a)g - T_a^{-1}(ag) = (a^{-1} a)g = g$ ).

	Observe next that $H=\{T_a: a \in G\} $ is a subgroup of $A(G ) $. For whenever $T_a, T_b \in H$, one has \[
		T_a \circ T_b^{-1} = T_{a} \circ T_{b^{-1} }=T_{ab^{-1} } \in H
	.\] 
	Hence $H \le A(G)$ by subgroup criterion.

	To justify the last inequality, note: if $a,b \in G$, for all $g\in G$, one has \[
		(T_a \circ T_b)(g) = T_a(bg) = abg=T_{ab }(g)
	,\]  
	so $T_a \circ T_b = T_{ab }$.

	Goal now: Show that $G \cong H \le  A(G)$.
	We need to construct an isomorphism $\varphi:G \to H$. Define \begin{align*}
		\varphi: G &\longrightarrow H \\
		a &\longmapsto \varphi(a) = T_a
	.\end{align*}
	This is a homomorphism, since for $a,b \in G$, one has  \[
		\varphi(ab) = \varphi(a)\varphi(b) =T_{ab }= T_a T_b = \varphi(a)\varphi(b)
	.\] 
	Injectivity: $\varphi(a) = \varphi(b) \iff T_a = T_b \iff \forall g \in G, ag=bg \iff a = b$.
	So $\varphi(a) = \varphi(b) \iff a = b$, so $\varphi$ is injective.

	Surjectivity: Whenever $T_a \in H$, one has $\varphi(a) = T_a$, so there exists $a \in G$ with $\varphi(a)=T_a$. So $\varphi$ is a bijective homomorphism (i.e. an isomorphism from $G$ to $H $). 
Thus $G \cong H \le A(G)$.
\end{proof}

\begin{lemma}
	(2.5.2 + 2.5.3) If $\varphi:G \to  G'$ is a homomorphism of groups then $\varphi(G)\le G'$.
\end{lemma}
\begin{proof}
	Homework exercise.
\end{proof}

\begin{definition}
	If $\varphi:G \to  G'$ is a homomorphism of groups, then the \textit{kernel} of $\varphi$ is defined by $\operatorname{ker} (\varphi) = \{a \in G: \varphi(a)=e'\}$.
\end{definition}

\begin{fact}
	$\operatorname{ker}\varphi \le G$.
\end{fact}

